\section{空间向量与解析几何}
\begin{enumerate}[itemsep=0.3em,topsep=0pt]
	
	\item 已知$M_1(4, \sqrt{2}, 1)$, $M_2(3,0,2)$, 求向量$\vec{M_1M_2}$的模、方向余弦、方向角.
	
	\item 在空间直角坐标系下，求 $\vec{a} = (4, -3, 4)$ 在 $\vec{b} = (2,2,1)$ 上的投影.

	\item $\vec{a} = 3\vec{i} - \vec{j} - 2\vec{k}$, $\vec{b} = \vec{i} + 2\vec{j} - \vec{k}$, 求 $2\vec{a} + 3\vec{b}$ 与 $5\vec{a} - 2\vec{b}$ 的内积和外积.

	\item 已知点$A(1,2,0)$, $B(2,3,1)$, $C(4,2,2)$, $M(x,y,3)$.\\
	(1) 若$A,B,M$三点共线, 求$x,y$的值;\\
	(2) 若$A,B,C,M$四点共面, 求$x,y$的关系.
	
	
	\item 求点$(1,1,-1)$, $(-2,-2,2)$, $(1,-1,2)$三点确定的平面方程.

	\item 求经过点$(2, -3, 1)$且与平面$3x - 7y + 5z = 12$平行的平面方程.

	\item 有平面$x - y + 2z = 1$和$2x + y + 4 = 0$.\\
	(1) 求两个平面的交线; \\ (2) 求包含这条交线，且过点$P(1,1,1)$的平面方程.

	\item 有直线$l_1: x-1=\frac{y+2}{2}=z-4$, $l_2: 2-x=\frac{y-3}{2}=z$, 求包含$l_2$，且与$l_1$平行的平面方程.

	
	\item 有直线$l_1: \frac{x+2}{2}=\frac{y}{-3}=\frac{z-1}{4}$, $l_2: \frac{x-3}{m}=\frac{y-1}{4}=\frac{z-7}{2}$, 当$m$取何值时，两条直线相交？并求出此时的交点.
	
	\item 求点$(2, -3, -1)$到直线$\frac{x-1}{-2}=\frac{y+1}{-1}=z$和平面$x+2y+2z-10=0$的距离和投影.

	\item 将$xOy$平面上的双曲线$4x^2 - 9y^2 = 36$分别绕$x$轴和$y$轴旋转一周，求所得旋转曲面的方程.

	\item 求曲线 $\begin{cases} y^2+z^2 = x \\ x+2y-z=0 \end{cases}$ 在$yOz$平面和$xOy$平面上的投影方程.
	
\end{enumerate}

\newpage

\section{多元函数微分学}
\begin{enumerate}
	\item 求极限： $\lim\limits_{(x,y) \to (0,0)} \frac{\ln(1+x^2+y^2)}{x^2+y^2}$ 

	\item 求极限： $\lim\limits_{(x,y) \to (0,0)} \frac{xy}{x^2+y^2}$

	\item 求极限： $\lim\limits_{(x,y) \to (0,0)} \left(x\sin\frac{1}{y} + y\sin\frac{1}{x}\right)$ 

	\item 求极限： $\lim\limits_{(x,y) \to (0,0)} (1+\sin x)^{\frac{y+1}{x}}$
	
	\item 判断函数 $f(x,y)=\begin{cases} 0 & x=y=0 \\ \frac{x^2y^2}{(x^2+y^2)^{\frac{3}{2}}} & \text{其他} \end{cases}$ 在$(0,0)$处是否连续.

	
	\item 设 $z=\frac{1}{2}\ln (x^2+y^2)+\arctan \frac{y}{x}$, 求 $\frac{\partial f}{\partial x}$, $\frac{\partial f}{\partial y}$.

	\item 求 $f(x,y) = x^2 \sin 2y$的二阶偏导 $\frac{\partial^2 f}{\partial x^2}$, $\frac{\partial^2 f}{\partial x \partial y}$, $\frac{\partial^2 f}{\partial y^2}$.

	\item 设 $z = z(x,y)$由方程$\frac{x}{z}=\ln \frac{z}{y}$确定, 求 $\frac{\partial z}{\partial x}$, $\frac{\partial z}{\partial y}$.

	\item 函数 $x = x(z)$, $y = y(z)$由方程组 $\begin{cases} x+y+z=0 \\ x^2+y^2+z^2=1 \end{cases}$ 确定, 求 $\frac{dx}{dz}$, $\frac{dy}{dz}$.

	\item 求函数 $z = xy + \frac{x}{y}$ 的全微分.

	
	\item 已知函数 $f(x,y)$的全微分为 $(x^2+axy)dx + (x^2+3y^2)dy$, 求$a$的值.

	\item 设 $z=xy^2+ye^{2x}$, 点P$(0,1)$ ,求:\\
	(1) $z$在P处沿$(-1,1)$方向的方向导数; \\ (2) $z$在P处方向导数的最大、最小值, 并求出相应方向.

	
	\item 求函数 $z = x^3 + y^2 - 6x^2 - 4y + 2$的极值.

	\item 求椭圆 $x^2 + 4y^2 = 4$ 上的点到直线 $2x+3y-6=0$ 距离的最小值.

	\item 求 $f(x,y) = 3x^2 + 3y^2 - x^3$ 在闭区域$D: x^2 + y^2 \leq 16$ 上的最值.

	\item 求曲面 $z = x^2(1-\sin y) + y^2(1-\sin x)$在$P(1,0,1)$处的切平面与法线方程.

	
	\item 求曲线$x=t-\sin t, y=1-\cos t, z=4\sin \frac{t}{2}$在$t=\frac{\pi}{2}$处的切线方程和法平面方程.

	\item 求曲线 $\begin{cases} x^2+y^2+z^2=6 \\ x^3+y+z^2=0 \end{cases}$ 在点 $P(1, -2, 1)$处的切线方程和法平面方程.

\end{enumerate}

\newpage

\section{多元函数积分学}
\begin{enumerate}
	\item 求重积分： $\iint\limits_D \sin(x+y)dxdy$, 其中$D$表示$0 \leq x \leq \pi$, $0 \leq y \leq \frac{\pi}{2}$ 的区域.

	\item 求重积分： $\iint\limits_D e^x \sin y dxdy$, 其中 $D$ 表示 $0 \leq x \leq 1$, $0 \leq y \leq \frac{\pi}{6}$ 的区域.	
	
	\item 求重积分： $\iint\limits_D (x+y)dxdy$, 其中 $D$ 表示 $y=x$ 和 $y=2-x^2$ 围成的区域.
	
	\item 求重积分： $\iiint\limits_D xyzdxdydz$, 其中 $D$ 表示满足 $0 \leq x \leq y \leq z \leq 1$ 的区域.
	
	\item 求重积分： $\iint\limits_D xy d\sigma$, 其中 $D$ 表示 $y^2 = x$ 和 $y=x-2$ 围成的区域.
	
	\item 求重积分： $\int_0^1 dy \int_y^{\sqrt{y}} \frac{\sin x}{x} dx$

	\item 求重积分： $\iint\limits_D xy^2 dxdy$, 其中 $D$ 表示 $y^2 = x$ 和 $y=x-2$ 围成的区域.
	
	\item 求重积分： $\iint\limits_D (x^3+\sin y+1)dxdy$, 其中 $D$ 表示 $x^2 + y^2 \leq r^2$ 的区域.
	
	\item 求重积分： $\iint\limits_D (x-2y)^2 dxdy$, 其中 $D$ 表示 $x^2 + y^2 \leq 4$ 的区域.
	
	\item 求曲面 $z = x^2 + 2y^2$ 与 $z = 6 - 2x^2 - y^2$ 围成部分的体积.
	
	
	\item 设$I_1=\iint\limits_{D_1} (x^2+y^2)d\sigma$, $I_2=\iint\limits_{D_2} 2|xy|d\sigma$, $I_3=\iint\limits_{D_2} (x^2+y^2)d\sigma$, 其中$D_1$表示$x^2+y^2 \leq 1$区域, $D_2$表示$|x|+|y| \leq 1$区域, 比较$I_1, I_2, I_3$的大小.
	
	\item 求曲线积分 $\int_L (x^2+y^2) ds$, 其中$L$表示半圆周 $y = -\sqrt{1-x^2}$.
	
	\item 求曲线积分 $\int_L y^2 ds$, 其中$L$表示$x = t - \sin t, y = 1 - \cos t \ (0 \leq t \leq 2\pi)$的部分.
	
	\item 求曲线积分 $\oint_L xds$, 其中曲线$L$表示$y=x^2$与$y=x$围成区域的边界.
	
	
	
	\item 设$L$表示椭圆$\frac{x^2}{4}+\frac{y^2}{3}=1$, 其周长为$a$, 求 $\oint_L (2xy+3x^2+4y^2) ds$.
	\item 求曲线积分 $\oint_L x^2 ds$, 其中$L$表示单位圆 $x^2 + y^2 = 1$.
	\item 设$L$表示$x^2 + y^2 = 2x$的正向, 求 $\oint_L y^2 dx + 2xdy$.
	\item 设$L$表示$y=x^2$与$x=y^2$围成曲线的正向, 求 $\oint_L (2xy-x^2)dx+(x+y^2)dy$.
	
	\item 设$L$表示曲线$x = t, y = t^2, z = t^3$上$0 \leq t \leq 1$的部分, 将 $\int_L Pdx+Qdy+Rdz$化为第一类曲线积分.
	
	\item 设$L$表示正向圆周$x^2+y^2=9$, 求 $\oint_L (2xy-2y)dx+(x^2-4x)dy$.
	
	\item 曲线$L$表示$x+y=1$在第一象限的线段与$x^2+y^2=1$在第二象限的圆弧所构成曲线的正向.
	求 $\int_L (x^2-2y)dx+(3x+ye^y)dy$
	
	\item 求 $\int_L (xy^2+2x)dx+(x^2y+2)dy$, 其中$L$表示一条从$(1,2)$到$(2,4)$的曲线.
	
	
	\item 验证$(xy^2+y)dx+(x^2y+x-2)dy$是某个函数$u$的全微分，并求出这个函数$u$.
	
	\item 记$S$表示锥面 $z=\sqrt{x^2+y^2}$ 与平面 $z=1$ 所围成区域的边界曲面.求曲面积分 $\iint\limits_S (x^2+y^2)dS$.
	
	\item 对于曲面 $\Sigma: z=\sqrt{4-x^2-y^2}$, 求 $\iint\limits_{\Sigma} (x+y+z)^2 dS$.

	\item 对于球面 $\Sigma: x^2+y^2+z^2=1$, 求 $\iint\limits_{\Sigma} (x^2+2y^2+3z^2)dS$.

	\item 对于球面 $\Sigma: x^2+y^2+z^2=1$ 的下半部分的下侧, 求 $\iint\limits_{\Sigma} (x^2+y^2+z)dxdy$.
	
	\item 设曲面$\Sigma$表示$z = x^2 + y^2 \ (0 \leq z \leq 4)$的上侧, 求 $\iint\limits_{\Sigma} xdydz + ydzdx + zdxdy$.
	
	\item 设$\Sigma$表示圆柱体$x^2+y^2 \leq 9 \ (0 \leq z \leq 3)$整个表面的外侧, 求 $\oiint\limits_{\Sigma} xdydz + ydzdx + zdxdy$.

	\item 设$\Sigma_1$表示柱面 $x^2+y^2 = 9 \ (0 \leq z \leq 3)$侧面的外侧, 求 $\iint\limits_{\Sigma_1} xdydz + ydzdx + zdxdy$.

	
	\item 设$L$是柱面$x^2+y^2=1$和平面$y+z=0$的交线, 从$z$轴正向往负向看去为逆时针方向, 求曲线积分 $\oint_L zdx + ydz$.
	
	\item 有向量场 $\vec{u}(x,y,z) = xy^2\vec{i} + ye^{z^2}\vec{j} + x\ln(1+z^2)\vec{k}$, 求向量场在 $P(1,1,0)$处的散度和旋度.
	
\end{enumerate}

\newpage

\section{无穷级数}
\begin{enumerate}
	
	\item 判断级数是否收敛
	\begin{enumerate}
		\item $\sum\limits_{n=1}^{\infty} \frac{1}{n(n-1)}$
		
		
		\item $\sum\limits_{n=1}^{\infty} \ln \left(1+\frac{1}{n}\right)$
		
		
		\item $\sum\limits_{n=1}^{\infty} \frac{1}{3n}$
		
		
		\item $\sum\limits_{n=1}^{\infty} \left(\frac{1}{2^n} + \frac{1}{3^n}\right)$
		
		
		\item $\sum\limits_{n=1}^{\infty} \sin \frac{\pi}{2^n}$
		
		
		\item $\sum\limits_{n=1}^{\infty} \frac{1+n}{1+n^2}$
		
		
		\item $\sum\limits_{n=1}^{\infty} \frac{1}{\int_0^n \sqrt[3]{1+x^3} dx}$
		
		\item $\sum\limits_{n=1}^{\infty} \frac{2^n n!}{n^n}$
		
		
		\item $\sum\limits_{n=1}^{\infty} \left(\frac{2n-1}{2n+1}\right)^{n^2}$
		
		
		\item $\sum\limits_{n=2}^{\infty} \frac{1}{n \ln n}$
		
		
		\item $\sum\limits_{n=2}^{\infty} \frac{1}{n(\ln n)^2}$
		
		
		\item $\sum\limits_{n=2}^{\infty} (-1)^n \frac{n+2}{n+1} \cdot \frac{1}{\sqrt{n}}$
		
		
		\item $\sum\limits_{n=1}^{\infty} \frac{(-1)^n \ln n}{\sqrt{n}}$
		
		
	\end{enumerate}
	
	\item 判断下列级数是否收敛，绝对收敛还是条件收敛：
	
	\begin{enumerate}
		\item $\sum\limits_{n=1}^{\infty} \frac{n^3 \sin \frac{n\pi}{3}}{2^n}$

		
		\item $\sum\limits_{n=1}^{\infty} \frac{(-1)^n}{n-\ln n}$

	\end{enumerate}
	
	\item 求下列级数的收敛域：
	\begin{enumerate}
		\item $\sum\limits_{n=1}^{\infty} \frac{(2x-1)^n}{n^2}$

		\item $\sum\limits_{n=1}^{\infty} \frac{3^n+(-2)^n}{n} (x-1)^n$

		\item $\sum\limits_{n=1}^{\infty} (-1)^n \frac{n}{2^n} (x-1)^{2n}$
	
	\end{enumerate}
	
	
	\item 将$f(x)=\frac{3x}{x^2+x-2}$ 展开为$x-2$的幂级数.

	
	\item 将$f(x)=x\ln(1-x-2x^2)$展开为$x$的幂级数.

	
	\item 将$f(x)=\arctan \frac{1+x}{1-x}$展开为$x$的幂级数.
		
	\item 将$f(x)=\frac{1}{(1+x)^2}$展开为$x-1$的幂级数.


	\item 求$\sum\limits_{n=1}^{\infty} \frac{x^{n+1}}{2^n \times n!}$的和函数.


	
	\item 求$\sum\limits_{n=1}^{\infty} \frac{x^n}{n(n+1)}$的和函数.
	
	\item 求$\sum\limits_{n=0}^{\infty} \frac{4n^2+4n+3}{2n+1} x^{2n}$的和函数.

	\item 求常数项级数$\sum\limits_{n=2}^{\infty} \frac{1}{(n^2-1) \times 2^n}$的和.

	\item 求常数项级数$\sum\limits_{n=1}^{\infty} \frac{1}{n! (n+2)}$的和.

	\item 将 $f(x)=x^2 \ (-\pi \le x < \pi)$展开为傅里叶级数.

	\item 有以$2\pi$为周期的函数$f(x)= \begin{cases} x+3 & 0 \le x \le \pi \\ 1-x & -\pi < x < 0 \end{cases}$, 设$S(x)$为$f(x)$的傅里叶级数的和函数，求$S(2\pi), S(\frac{5\pi}{2}), S(9\pi)$.


	
\end{enumerate}

\newpage

\section{常微分方程} 
\begin{enumerate}
	\item 求方程 $y^{\prime}=\sqrt{y}$ 的通解.

	\item 解初值问题$ \left\{\begin{array}{l} x y \mathrm{~d} x+\left(x^{2}+1\right) \mathrm{d} y=0 \\ y(0)=1 \end{array}\right.$
	
	\item 求解微分方程$\frac{d y}{d x}=2 x y$的通解.

	
	\item 解微分方程 $y^{\prime}=\frac{y}{x}+\tan \frac{y}{x}$.

	
	\item 求方程$ y^{\prime}=\frac{y}{x+y^{3}} $的通解.

	\item 求方程$ \frac{d y}{d x}-\frac{4}{x} y=x^{2} \sqrt{y} $的通解.

	
	\item 求解$ y^{\prime\prime\prime}=e^{2x}-\cos x$.

	\item 求解$ \left\{ \begin{array}{l} \left( 1+x^2 \right) y^{\prime\prime}=2xy^{\prime}\\ \left. y \right|_{x=0}=1,\left. \quad y^{\prime} \right|_{x=0}=3\\ \end{array} \right.$

	
	\item 求$ y^{\prime\prime\prime}=y^{\prime\prime} $的通解.

	
	\item 求解$\frac{dy}{dx} = \frac{x^2+2}{y}$.

	
	\item 求解$\frac{dy}{dx} + 6x^2y = e^{-2x^3}\sin(x)$.

	\item 求$y''-6y'+9y=4e^{3x}$的通解.

	\item 求$y''-7y'+12y=5$的通解.
	
	\item 求$y''-2y'+4y=(x+2)e^{3x}$的通解

	\item 求$y''+2y'+5y=\sin 2x$的通解.

	\item 求$y''-2y'+10y=x\cos 2x$的通解.

	
\end{enumerate}
